\chapter{金属键：为什么铜能拉成电线}

\begin{kaichang}
找一根废旧的电线，请大人帮你剥开外面的塑料皮（只看，不要通电）。里面是一束暗红色的细丝——铜丝。试试用手弯它：弯得动，而且弯过去就待在那里，不断也不弹回来。再找一根差不多粗的玻璃棒或者一根粉笔，也弯一弯：啪，断了。

再来看看这根铜丝干的活儿。它被工厂从一大块铜坯拉成比头发丝还细的线，拉了几千米长也没有断成几截；它被裹进塑料皮里通上电，能让电灯亮、让手机充电。同样是一堆原子堆出来的固体，为什么铜可以被无限“拉丝”，还能让电流长驱直入，而玻璃和粉笔一弯就断、一碰电就“装死”？

先别急着说“因为铜是金属”。这句话只是把问题换了个名字。真正的问题是：金属里面的原子，到底是用一种什么样的方式“抱团”的？这一章，我们就把这根铜丝一路放大，放大到原子那一层去看。
\end{kaichang}

\section{金属们的四张共同面孔}

先把全世界八十多种金属摆在一起——铁、铜、铝、金、银、锡、锌、镍、钛……它们颜色不同、软硬不同、贵贱悬殊，但几乎都有四张共同的面孔。

第一张脸：\textbf{导电}。家里的电线用铜，高压输电线用铝，插座里的弹片用铜合金——凡是需要电流走的路，几乎都交给金属。而塑料、陶瓷、玻璃、干燥的木头，都被用来“拦住”电流。

第二张脸：\textbf{导热}。炒锅用铁或铝做，就是因为锅底上的热要飞快传给菜；而锅把手一定换成木头或塑料，不然你的手就是菜的一部分。冬天摸同一张桌子，铁桌腿总比木桌面凉——这个细节我们先记下来，章末有大用。

第三张脸：\textbf{能变形而不碎}。铜能拉成细丝，铝能压成比纸还薄的箔，金更夸张：一克金子可以锤成一张接近一平方米的金箔，薄到透光。工匠管这种本事叫“延展性”。你反复弯一根回形针，它会越弯越硬、最后才断；而食盐晶体、玻璃、陶瓷，稍微一掰就裂——这两种“死法”完全不同，本章后面你会看到这个对比有多关键。

第四张脸：\textbf{金属光泽}。新切开的金属总是亮闪闪的，金、银、铂靠这张脸当了几千年首饰。磨成粉的铝和银刷在漆里就是“银粉漆”，镀在玻璃背面就是镜子。

四种性格，出现在八十多种脾气各异的元素身上——第 10 章说过，周期表里同一类性格背后必有共同原因。铜、铁、铝又不会开会，它们凭什么整齐划一？线索要回到电子身上。

\begin{shiyan}
\textbf{材料}：一节 9 伏电池（或两节 5 号电池串联）、一颗发光二极管（LED）或手电筒小灯泡、几根导线；待测物品：一把金属钥匙、一枚硬币、一根铅笔芯（从铅笔里取出或直接用自动铅笔芯）、一把塑料尺、一块橡皮、一张折叠的铝箔、一个回形针。

\textbf{步骤}：
\begin{enumerate}[itemsep=1pt]
\item 把电池、LED 和导线连成一个故意“断开一口”的回路（LED 长腿接电池正极方向）。
\item 用待测物品挨个去“补”这个缺口：物品两端分别接触断开的两个线头，看 LED 亮不亮，并把结果记成两列：亮 / 不亮。
\item 再做一个“弯折测试”：把回形针掰直再弯折几次，感受它的变化；再取一支粉笔，轻轻弯一弯，观察发生了什么。
\end{enumerate}

\textbf{观察点}：哪些物品能让 LED 亮？铅笔芯居然也导电——它不是金属，而是石墨，第 24 章会专门讲这位“非金属里的特例”。回形针弯折几次后是不是变烫、变硬、最后才断？粉笔的“死法”和它有何不同？

\textbf{安全提示}：本实验只许使用电池，电压不超过 9 伏；\textbf{绝对禁止}用插座里的市电做任何尝试；LED 不要直接长时间接在电池两端（中间要串接待测物或一个小电阻），以免烧坏。
\end{shiyan}

\section{把镜头推进：一座泡在电子海里的城市}

回忆第 10 章和第 11 章的结论：金属元素集中在周期表左半边和中间，它们的共同点是\textbf{最外层电子少，而且离核较远、束缚较松}。钠最外层只有 1 个电子，镁 2 个，铝 3 个，连铁、铜这些过渡金属，最外层的电子也“待得不那么安稳”。

现在想象亿万个铜原子挤在一起。每个铜原子最外层的那颗电子，受到自家原子核的吸引力本来就不算强，旁边又挤满了别家原子核——对它来说，“留在自己家”和“溜达到邻居家”在能量上几乎没有区别。于是这些外层电子干脆\textbf{不再属于任何一个特定的原子}，而是变成整个金属块共有的“公共财产”，可以在整块金属里自由游荡。

剩下的部分呢？丢了外层电子的金属原子变成带正电的\textbf{阳离子}，它们排成整齐的阵列，像一座规划严整的城市里一栋栋位置固定的楼房；而那些自由电子像海水一样，把整座城市泡在里面。

\begin{center}
\begin{tikzpicture}
\foreach \x in {0,1.4,2.8,4.2} {
  \foreach \y in {0,1.4,2.8} {
    \draw[fill=gray!25] (\x,\y) circle (0.38);
    \node at (\x,\y) {$+$};
  }
}
\fill (0.7,0.5) circle (0.05);
\fill (2.1,0.9) circle (0.05);
\fill (3.5,0.4) circle (0.05);
\fill (0.3,1.1) circle (0.05);
\fill (1.5,1.9) circle (0.05);
\fill (2.9,2.2) circle (0.05);
\fill (3.9,1.6) circle (0.05);
\fill (0.9,2.6) circle (0.05);
\draw[->,thick] (4.6,0.2) -- (5.4,0.2) node[right]{\footnotesize 电子可以自由移动};
\end{tikzpicture}

\footnotesize 金属内部的人为表示法：大圆是位置固定的金属阳离子，小黑点是到处游荡的离域电子。圆圈、比例和位置都不代表真实外观。
\end{center}

这就是\textbf{金属键}：不是一根根把原子两两拴住的小棍（第 17 章拆过这个比喻），而是“整齐排列的阳离子”与“四处流动的电子海”之间的电性吸引。电子带负电，阳离子带正电，异号相吸，把整块金属粘成一个能量更低的整体——成键放能、断键吸能这条总账，和第 17 章讲过的任何化学键一样。

拿它和前两章的邻居比一比，差别就清楚了。离子键（第 18 章）里，电子被整个\textbf{过户}给氯这样的原子，从此死守在新主人身边；共价键（第 19 章）里，电子被两个原子\textbf{合伙持有}，活动范围就是两个核之间那一小片；而金属键里，电子是\textbf{全员公有}，活动范围是整块金属。化学家管这种“电子不归某两个原子私有、属于整个体系”的状态叫\textbf{离域}。三种键，本质都是电荷之间的吸引，区别只在电子的“户籍制度”：私有、合伙，还是全民所有。

\section{一副模型，四张脸全解释}

“阳离子泡在电子海里”这副图景，威力大得惊人——金属的四张共同面孔，它能一口气全解释。

\textbf{导电}：电子海里全是能自由移动的带电粒子。平时它们像广场上游荡的人群，各个方向乱走，平均下来哪也没去。一旦在金属两端接上电池、加上电场，人群虽然还在乱走，但整体多了一个缓慢的定向漂移——这就是电流。塑料和玻璃里没有这样的自由电子：它们的电子都被共价键“拴死”在特定的原子对之间，电场来了也无兵可派。

\textbf{导热}：你在锅底点火，那头的金属原子振动加剧，自由电子跑过去一撞，能量就交给了电子；电子跑得飞快，转眼撞到远处的离子，能量就传开了。电子海既是电流的高速公路，也是热量的高速公路——难怪导电好的金属导热也好，铜、铝、银两头都拔尖。

\textbf{延展}：这是金属键最得意的一场戏。锤打金属时，整齐排列的阳离子层会彼此滑动错开。要是换成离子晶体（回想第 18 章的食盐），层一滑动，原来正负相间的邻居就错成了正对着正、负对着负，同性相斥，晶体当场崩裂——所以盐和玻璃一掰就碎。金属却不怕：阳离子层滑过去，电子海立刻漫过来填上，负电的“海水”照旧把新位置的正离子粘住。金属键\textbf{没有方向}，不挑邻居的位置，怎么滑都不断——这就是铜能拉成几千米细丝的原因。

\textbf{光泽}：自由电子还有个本事：它们能吸收照进来的各种颜色的光，又几乎原样把光重新“吐”出去。于是金属表面像一面由电子组成的镜子，亮闪闪的。铜偏红、金偏黄，是因为它们对不同颜色的光反射效率略有差别——这背后已经是更深一层的量子规则了。

\begin{qiaoliang}
“自由电子”到底有多少？走得多快？我们来算一笔账。铜的密度约 \ud{8.9}{g/cm^3}，摩尔质量约 \ud{64}{g/mol}，所以 \ud{1}{cm^3} 铜里约有 $8.9\div 64\approx 0.14$ 摩尔铜原子，即 $0.14\times 6.02\times 10^{23}\approx 8.4\times 10^{22}$ 个原子。每个铜原子交出 1 个外层电子入海，那么一粒方糖大小的铜块里，自由电子就有 $8.4\times 10^{22}$ 个——比地球上的人口多十万亿倍。

再看看电流里它们走多快。假设截面积 \ud{1}{mm^2} 的铜导线里通着 \ud{1}{A} 的电流（大致是一台小台灯的电流），用电流等于“电子密度 × 电荷 × 漂移速度 × 截面积”可以算出，电子定向漂移的速度大约只有 $7\times 10^{-5}$ \ud{}{m/s}，也就是\textbf{每秒约 0.07 毫米}——比蜗牛爬还慢几十倍！可你按开关的瞬间灯就亮，这是为什么？因为“出发”的信号不是某个电子从开关一路跑到灯泡，而是电场以接近光速的速度传遍整条导线，命令导线里\textbf{本来就到处都有}的自由电子同时起步漂移。电线从不需要等电子“到货”，它是一座早就住满电子的城市。
\end{qiaoliang}

\section{符号登场：不写“分子式”的物质}

到这里，该让符号出场了。铜 \chem{Cu}、铁 \chem{Fe}、铝 \chem{Al}、银 \chem{Ag}、金 \chem{Au}——注意到没有，金属的“化学式”就是元素符号本身。这不是偷懒，而是在陈述一个事实：一块铜里没有任何可以圈出来的“铜分子”，它是亿万个 \chem{Cu} 原子（准确说是 \chem{Cu} 阳离子加电子海）组成的巨大整体，最小的“结构单元”就是原子。这和食盐一样：第 18 章说过，\chem{NaCl} 不表示一个个孤立的“食盐分子”，只表示晶体里钠离子和氯离子的比例。金属连“比例”都谈不上——满城只有一种居民。

也正因为金属里没有分子，说“一个铜分子”和说“一个食盐分子”犯的是同一种错误。这一细节在做化学计算时会很顺手：\ud{64}{g} 铜就是 1 摩尔铜原子，没有中间商。

\section{科学家怎样知道：电子海不是拍脑袋想的}

\begin{zhengju}
“金属里有一片电子海”听起来像个聪明的故事，凭什么相信它？证据来自一次漂亮的“跨界对账”。

1853 年，德国物理学家维德曼和夫兰兹测量了一大堆金属的两项性能：导电能力和导热能力。结果发现一个惊人的规律：在同一温度下，各种金属的\textbf{热导率除以电导率，得到的数值几乎一样}，而且这个比值随温度升高成正比地变大。银、铜、铁、铅，脾气差那么多，这项比值却像约好了似的。这意味着什么？意味着导电和导热很可能是\textbf{同一群载体}干的活——如果是两群互不相干的粒子分别管电和热，没有理由出现这么整齐的比值。

1900 年，也就是电子被发现（第 7 章提过汤姆孙的实验）仅仅三年后，德鲁德正式提出：金属内部有一团“电子气”，像气体分子一样在阳离子的缝隙里自由穿行。用这个模型一算，维德曼—夫兰兹的比值居然能对上号——电子海模型拿到了第一份硬证据。

但故事没有停在胜利。德鲁德的模型很快撞上南墙：按当时的统计物理，如果每个原子真贡献一个自由电子，这些电子就应该像气体分子一样吸收热量，让金属的“储热能力”（热容）明显大于非金属固体。可实验测出来，金属和非金属固体的热容几乎一样——那几万亿亿个“自由电子”在吸热这件事上仿佛不存在。模型解释得了导电，却解释不了热容，一定有什么地方错了。

修正要等到量子力学。1927—1928 年，索末菲和布洛赫等人把电子当作量子世界里的“波”重新计算，答案水落石出：电子海里的电子不是人人平等的自由身，它们必须按能量从低到高“入座”，绝大多数座位周围已经满员，动弹不得；只有能量最高、紧挨着“海面”（物理学家叫它费米面）的一小撮电子，才有空位可去、能真正参与导电导热。算一算比例：室温下能自由行动的电子大约只有总数的百分之一量级——难怪它们对热容的贡献小到几乎测不出。矛盾解决了，电子海模型升级成了更精细的版本。你看，科学模型的成长方式不是“提出即真理”，而是先解释一批事实，再被新事实逼到墙角，然后长出新的牙齿。
\end{zhengju}

\section{从电子海到能带：为什么是铜导电、硫不导电}

量子修正还顺手回答了一个更根本的问题：凭什么金属有自由电子，非金属就没有？

想法是这样的。一个孤立原子的电子住在分立的“楼层”上（第 9 章）。现在把 $N$ 个原子挤到一起，每个原子的楼层都会受到邻居的影响，一层楼就“劈裂”成 $N$ 条靠得极近的细线。$N$ 是一亿亿亿量级，细线密到连成一片——这一片就叫一条\textbf{能带}。带与带之间可能空着一段没有任何楼层的禁区，叫\textbf{禁带}。

\begin{itemize}[itemsep=2pt]
\item \textbf{金属}（铜、铁、铝）：最外面那条能带恰好没住满，或者和上面一条空带重叠。电子稍微加点能量，旁边就有空位可搬——于是它们能像电子海里描述的那样自由行动，导电导热。
\item \textbf{绝缘体}（硫、玻璃、金刚石）：最外面一条带住得满满当当，再往上隔着一条很宽的禁带。满员车厢里谁也挪不动窝，普通电压也提供不了跨越禁带的能量——电子被“焊死”，不导电。
\item \textbf{半导体}（硅、锗）：禁带很窄，热运动或掺杂能让少量电子跨过去——导电能力介于两者之间，而且可以被人为精确调控。你的手机芯片就是靠这种“可调”的脾气工作的（第 10 章提过硅的这个身份）。
\end{itemize}

所以电子海模型和能带理论是同一座大厦的两层：前者告诉你“自由电子干了什么”，后者告诉你“自由电子从哪里来、为什么有的材料一个都没有”。对初中生来说，电子海够用；能带这扇门，先推开一条缝，看到里面有光就行。

\section{往城里搬几户“外乡人”：合金}

纯金属其实有个职场弱点：正因为离子层容易滑动，它们普遍偏软。纯铁软得能用手掰弯，纯铝一按一个坑，纯金软到不能用做戒指。而人类几千年来一直在用同一个办法给金属“上强度”：\textbf{掺别的原子，做成合金}。

道理就在金属键里。层与层之所以容易滑动，是因为每一层都整整齐齐。现在搬几户个头不同的“外乡人”进来：大个子的外来原子挤在队伍里，小个子的原子钻进队伍的空隙——原本平滑的层间通道变得坑坑洼洼，再想滑动就费劲了。滑动难，材料就硬、就强。电子海照样把大家粘在一起，所以合金依然保留导电、有光泽这些金属本色。

这张“掺”出来的菜单，你多半吃过、用过、坐过：

\begin{itemize}[itemsep=2pt]
\item \textbf{钢}：铁里掺进少量碳（质量分数通常在 2\,\% 以下）。碳原子很小，正好钻进铁原子阵列的空隙里，把层间通道“卡住”。铁就此从软汉变成硬汉——桥梁、铁轨、钢筋全靠这一掺。碳掺多了又变脆，所以炼钢的火候本质是控制“卡子”的数量。再掺入铬（约 12\,\% 以上）和镍，表面会形成一层极薄的氧化铬保护膜，锈就很难落脚——这是不锈钢（铬的这种本领，第 15 章逛过渡金属街区时还会遇到）。
\item \textbf{青铜}：铜掺锡。这是人类学会的第一种合金，三千多年前就把历史推进了“青铜时代”——它比纯铜硬得多，熔点还更低，更好浇铸，鼎、剑、钟都是它。
\item \textbf{黄铜}：铜掺锌。金黄色的它常冒充黄金，做乐器、水龙头和弹壳——锌原子比铜略大，替换进铜的队伍里，同样起到“卡位”效果。
\item \textbf{铝合金}：铝最大的优点是轻（密度只有铁的三分之一），缺点是软。掺入少量铜、镁、硅之后，强度能赶上钢，体重却轻得多——飞机机身、易拉罐、自行车架都是这笔买卖。
\item \textbf{形状记忆合金}：镍和钛按约一比一的原子比例组成的合金有个魔术：把它掰弯，再用温水一浇，它自己弹回原来的形状。原理藏在晶体层面——它的原子阵列能在两种排列方式之间随温度切换，像一支记住了两种队形的仪仗队。眼镜架、矫正牙齿的弓丝、心脏血管支架，都用上了这种“记得自己是谁”的金属。
\end{itemize}

合金不是化合物。铁和碳、铜和锌之间没有固定的组成比例，也写不出化学式——它们更像一锅调好比例的杂烩，原子们在晶格里混居，靠的还是那片电子海。也正因如此，合金的性能可以像调配方一样连续调节，这是材料工程师的看家本领。

\section{这副模型在哪里会撞墙}

每副模型都有它的边界，电子海也不例外。它擅长解释\textbf{共性}——为什么金属都导电、都导热、都能延展、都有光泽。但一碰到\textbf{个性}，它就开始含糊：

\begin{itemize}[itemsep=2pt]
\item 为什么铁、钴、镍能被磁铁吸住，铜和铝却不能？磁性来自电子自旋的集体排队，电子海模型里根本没有这条。
\item 为什么钨的熔点高达 $3422\,^{\circ}\mathrm{C}$，而汞在 $-39\,^{\circ}\mathrm{C}$ 就熔化、常温下是液体？电子海说大家都是“阳离子泡在海里”，讲不清这近三千五百度 的差距。
\item 为什么有些金属冷到接近绝对零度时，电阻会突然消失得无影无踪，电流可以在里面永远流下去（超导）？这需要电子两两结伴的量子图景，1911 年昂内斯发现它时，整个物理学界都懵了——直到今天，高温超导的完整解释仍是未完全解决的难题。
\end{itemize}

这些墙后面站着的都是量子力学。所以请这样安放本章的模型：定性解释金属的四大共性，它是最好的第一副眼镜；要预测某种具体金属的硬度、熔点、磁性，就得换更精细的镜片。

\begin{wuqu}
“金属摸起来凉，是因为金属的温度比周围低。”——错得很有迷惑性。把一根铁条和一根木条在同一间屋里放一晚上，它们的温度必然相同（都等于室温），温度计一量便知。差别在于：你的手指大约 $33\,^{\circ}\mathrm{C}$，比室温高。摸木头时，木头导热慢，指尖的热量散不掉，皮肤表面很快暖起来，你觉得温；摸铁条时，电子海像海绵一样把你指尖的热量迅速抽走、摊到整条铁条里，皮肤表面一直暖不起来，于是你觉得凉。“凉”不是铁的温度低，而是你的热量流失快。同一个道理反过来用：刚出炉的铁锅比同温度的木铲烫得多——因为它往你手里灌热量的速度也快得多。导热性描述的是\textbf{热量迁移的速率}，不是物体的冷热，速度和多少永远是两回事。
\end{wuqu}

\section{回到那根铜丝}

现在，重新剥开那根电线。

铜丝能被拉成几千米长而不断，因为铜 \chem{Cu} 的阳离子层在拉拔时层层滑移，而离域电子汇成的海立刻漫过每一个新位置，粘住全城——金属键没有方向，滑到哪里，哪里还是家。玻璃和粉笔一弯就断，因为它们的原子靠有固定方向的共价键、或正负相间的离子键连接，位置一错，吸引力就变成了排斥力。

铜丝能通电，因为每个铜原子交出的那颗外层电子不归任何原子私有，亿亿亿亿个自由电子本来就挤满了导线全程；开关一合，电场以近光速发令，全城电子同时开始每秒不到一毫米的缓慢漂移——灯立刻亮，不需要任何电子从发电厂“快递”到你家。

电线外面那层塑料皮呢？它的分子里，电子被共价键牢牢拴在特定的原子对之间（第 19 章），满员车厢，禁带又宽，一个自由电子也派不出来——所以它是绝缘体，保护你不被电流误伤。一根电线，就是“电子海”和“电子私有制”并肩工作的现场：里面的铜负责让电子流动，外面的塑料负责让电子别乱流动。你看，化学键的类型直接决定了材料的分工。

\section{下一站的伏笔}

这一章我们看到，铜块内部靠“阳离子泡在电子海里”凝成整体；第 18 章的食盐靠正负离子的吸引排成晶体；第 19 章的水分子靠共享电子在分子内部紧紧抱团。

可是，这里有一个被我们一直跳过的问题。水分子内部的氢和氧靠共价键结合，那\textbf{一个水分子和另一个水分子之间}呢？水能聚成一滴、糖能溶在水里、壁虎能倒挂在天花板上、保鲜膜能自己粘在碗上——这些分子和分子之间，既不够成键，又确实存在某种“弱弱的牵扯”。金属键、离子键、共价键都管不了这件事，因为它们只管“原子怎样结成分子或晶体”，不管“分子之间怎样相处”。这种更微弱、却同样塑造了整个世界的吸引力是什么？第 22 章，我们就去抓这股藏在分子之间的手。

\begin{tiaozhan}
为什么满员能带里的电子“动弹不得”，而费米面附近的电子可以？量子力学里，电子是波，每个电子必须占据一个独一无二的量子状态（泡利不相容原理，第 9 章挑战层见过面）。要让一个电子动起来参与导电，它就得加速、改变运动状态，也就是“搬家”到另一个能量稍高的空状态去。满带里所有邻近座位都住了人，电子无处可搬，只能全体原地待命；而金属的费米面附近，满座与空座紧挨着，加一点点能量（电场的推力）就有电子搬得动家。进一步可以算出：热运动只能“唤醒”能量在费米面上下约 $kT$ 范围内的电子，$kT$ 相对于费米能只有百分之一量级——这正是德鲁德热容之谜的定量答案：真正自由的电子，从来都只是“海面上的一层浪花”。这部分跳过不影响主线，但如果你将来想知道芯片、超导和激光为什么能存在，答案全都埋在这条费米面附近。
\end{tiaozhan}

\section*{练一练}

\begin{sikaoti}
\item 画一幅金属内部的粒子示意图：用带“$+$”号的大圆表示整齐排列的阳离子，用小点表示离域电子。再画一幅“层被推拉后错开”的图，标出电子海怎样填上新的空隙。图旁注明：这是帮助思考的表示法，不代表真实外观。
\item 用电子海模型逐项解释：为什么（1）铝能压成箔、（2）铜线能导电、（3）银勺放进热汤里勺柄很快变烫。每一项分别对应模型里的哪个角色（阳离子阵列还是自由电子）？
\item 一根电线断了，有人用胶带把两端铜丝拧在一起缠好。从“电子怎样流过接头”的角度想一想：为什么拧得越紧、接触面越大，接头处越不容易发热？（提示：电流喜欢走宽敞的路。）
\item 把一把金属勺和一把塑料勺同时放进冰箱冷冻室，一小时后同时拿出，用手背分别碰一下，哪把更凉？它们的温度真的一样吗？用本章的模型解释你的手背“感觉”到了什么。
\item 纯金很软，做成首饰容易刮花变形。金匠掺入少量银或铜做成“K 金”，硬度大增。画出纯金晶体中一层原子滑动的示意图，再画出掺入大小不同的原子后滑动受阻的示意图，说明合金为什么更硬。
\item 钠最外层有 1 个电子，镁有 2 个，铝有 3 个。猜测：这三种金属里，谁对电子海的“贡献”最大、金属键可能最强、熔点可能最高？查一查三种金属的熔点，验证你的猜想。
\end{sikaoti}
